\clearpage
\pdfbookmark[0]{Ergänzung zur Theorie der Kombination von Beobachtungen}{supplementum}
\thispagestyle{empty}
\SourceLeaf{060}
\begin{center}
\vspace*{2em}
{\Large ERGÄNZUNG\par}
\bigskip
{\Large DER THEORIE\par}
\bigskip
{\Large\bfseries DER KOMBINATION VON BEOBACHTUNGEN\par}
\bigskip
{\large DIE DEN KLEINSTEN FEHLERN AUSGESETZT IST\par}
\bigskip
VON\par
\bigskip
{\large CARL FRIEDRICH GAUSS\par}
\bigskip
{\small DER KÖNIGLICHEN GESELLSCHAFT DER WISSENSCHAFTEN IM JAHR 1826, AM 16. SEPTEMBER VORGELEGT.\par}
\bigskip
\rule{.7\linewidth}{.3pt}\par
\bigskip
Commentationes societatis regiae scientiarum Gottingensis recentiores.
\mbox{Vol.~\textsc{vi}.}
\par\medskip
Göttingen MDCCCXXVIII.
\par\bigskip
\rule{.7\linewidth}{.3pt}
\end{center}
\clearpage
\renewcommand{\Article}[1]{\par\medskip\Needspace{4\baselineskip}\pdfbookmark[1]{Art. #1}{supplementum-art-#1}\centerline{\textbf{#1.}}\nobreak\smallskip}
\SourceLeaf{062}
\begin{center}
{\large ERGÄNZUNG\par}
\medskip DER THEORIE\par
\medskip {\large DER KOMBINATION VON BEOBACHTUNGEN\par}
\medskip DIE DEN KLEINSTEN FEHLERN AUSGESETZT IST.\par
\medskip\rule{.25\linewidth}{.3pt}
\end{center}
\Article{1}
In der in Band V der Commentationes Recentiores enthaltenen Behandlung der Theorie der Kombination von Beobachtungen haben wir vorausgesetzt, dass jene Größen, deren Werte durch Beobachtungen ohne absolute Genauigkeit gegeben sind, so von gewissen unbekannten Elementen abhängen, dass sie als gegebene Funktionen dieser Elemente dargestellt sind, und dass es im Wesentlichen darauf ankommt, diese Elemente möglichst genau aus den Beobachtungen abzuleiten.

In den meisten Fällen trifft diese Voraussetzung tatsächlich unmittelbar zu. In anderen Fällen stellt sich die Aufgabe jedoch in einer etwas anderen Form dar, sodass es auf den ersten Blick zweifelhaft erscheint, wie sie auf die verlangte Form zurückgeführt werden kann. Nicht selten kommt es nämlich vor, dass die Größen, auf die sich die Beobachtungen beziehen, noch nicht als Funktionen bestimmter Elemente dargestellt sind und auch nicht auf eine solche Form zurückführbar erscheinen, zumindest nicht bequem oder ohne Umwege; während andererseits die Natur der Sache gewisse Bedingungen vorgibt, denen die wahren Werte der beobachteten Größen genau und notwendig genügen müssen.

Bei näherer Betrachtung lässt sich jedoch leicht erkennen, dass sich dieser Fall vom anderen tatsächlich nicht wesentlich unterscheidet, sondern auf ihn zurückgeführt werden kann. Bezeichnen wir nämlich die Anzahl der beobachteten Größen mit \(\pi\), die Anzahl der Bedingungsgleichungen aber mit \(\sigma\) und wählen wir von den ersteren beliebige \(\pi-\sigma\) aus, so hindert uns nichts daran, diese selbst als die Elemente anzunehmen und die übrigen, deren An\SourceLeaf{063}%
zahl \(\sigma\) sein wird, mithilfe der Bedingungsgleichungen als Funktionen jener zu betrachten; so ist die Sache auf unsere Voraussetzung zurückgeführt.

Doch auch wenn dieser Weg in sehr vielen Fällen recht bequem zum angestrebten Ziel führt, lässt sich nicht leugnen, dass er weniger natürlich ist und es sich daher lohnt, die Aufgabe in dieser anderen Form gesondert zu behandeln, umso mehr, als sie eine sehr elegante Lösung zulässt. Da diese neue Lösung sogar zu rascheren Rechnungen führt als die Lösung der Aufgabe in ihrer früheren Form, sobald \(\sigma\) kleiner als \(\frac12\pi\) ist oder, was dasselbe ist, sobald die in der vorhergehenden Abhandlung mit \(\rho\) bezeichnete Anzahl der Elemente größer als \(\frac12\pi\) ist, wird es in einem solchen Fall zweckmäßig sein, die neue Lösung, die wir in der vorliegenden Abhandlung erläutern werden, der früheren vorzuziehen, sofern sich die Bedingungsgleichungen ohne Umwege aus der Natur der Aufgabe gewinnen lassen.

\Article{2}
Bezeichnen wir mit \(v,v',v''\) etc. die \(\pi\) Größen, deren Werte durch Beobachtung bekannt werden, und hänge eine unbekannte Größe so von ihnen ab, dass sie durch eine gegebene Funktion derselben, etwa \(u\), dargestellt wird; ferner seien \(l,l',l''\) etc. die Werte der Ableitungen
\[
\frac{\mathrm du}{\mathrm dv},\quad \frac{\mathrm du}{\mathrm dv'},\quad \frac{\mathrm du}{\mathrm dv''}\ \text{etc.}
\]
die den wahren Werten der Größen \(v,v',v''\) etc. entsprechen. Wie nun durch Einsetzen dieser wahren Werte in die Funktion \(u\) deren wahrer Wert hervorgeht, so erhält man, wenn man für \(v.\ v',v''\) etc. Werte einsetzt, die von den wahren jeweils um die Fehler \(e,e',e''\) etc. abweichen, einen fehlerhaften Wert der Unbekannten, dessen Fehler sich schreiben lässt als
\[
=le+l'e'+l''e''+\text{etc.}
\]
sofern, wie wir stets voraussetzen werden, die Fehler \(e,e',e''\) etc. so klein sind, dass man (bei einer nicht linearen Funktion \(u\)) ihre Quadrate und Produkte vernachlässigen darf. Obwohl die Größe der Fehler \(e,e',e''\) etc. ungewiss bleibt, lässt sich die einer solchen Bestimmung der Unbekannten anhaftende Unsicherheit dennoch allgemein abschätzen, und zwar durch den bei einer solchen Bestimmung zu befürchtenden mittleren Fehler, der nach den Grundsätzen der vorhergehenden Abhandlung gleich dem folgenden Ausdruck ist:
\[
=\BareRoot{(llmm+l'l'm'm'+l''l''m''m''+\text{etc.})}
\]
\SourceLeaf{064}%
wobei \(m,m',m''\) etc. die mittleren Fehler der Beobachtungen bezeichnen, oder, wenn die einzelnen Beobachtungen gleicher Unsicherheit unterliegen,
\[
=m\BareRoot{(ll+l'l'+l''l''+\text{etc.})}
\]
Offenbar dürfen in dieser Rechnung für \(l,l',l''\) etc. mit gleichem Recht auch jene Werte der Ableitungen angenommen werden, die den beobachteten Werten der Größen \(v,v',v''\) etc. entsprechen.

\clearpage
